\section{Domain Task 3: Phân tích Cơ cấu Nguồn cung theo Khu vực}

Câu hỏi nghiệp vụ: ``Cơ cấu loại phòng theo từng quận phân bố như thế nào,
và khu vực nào phù hợp với từng quy mô nhóm khách?''

Mục đích: Giúp hiểu rõ đặc trưng phân khúc thị trường của từng khu vực và
nhu cầu sức chứa để định vị sản phẩm phù hợp.

\subsection{Biểu đồ 1: Cơ cấu loại phòng theo quận (Stacked Bar Chart)}

\subsubsection*{A. Thiết kế Idiom}

\begin{longtable}{@{} p{2.8cm} p{12cm} @{}}
\toprule
\textbf{Đặc điểm} & \textbf{Chi tiết} \\
\midrule
\endhead

\bottomrule
\endfoot

\textbf{Idiom} & Stacked Bar Chart \\
\addlinespace[0.2cm]

\textbf{What} & Khu vực: Categorical, Loại phòng: Categorical, Số lượng listing: Quantitative \\
\addlinespace[0.2cm]

\textbf{Why} & produce $\rightarrow$ browse, compare $\rightarrow$ summarize
\begin{itemize}[nosep, leftmargin=*, topsep=0.1cm]
    \item Biểu đồ giúp so sánh cơ cấu nguồn cung Airbnb giữa các khu vực, đồng thời cho biết loại phòng nào chiếm ưu thế trong từng khu vực.
\end{itemize} \\
\addlinespace[0.2cm]

\textbf{How} & \textbf{Encode:}
\begin{itemize}[nosep, leftmargin=*, topsep=0.1cm]
    \item Mark: Bar
    \item Channel:
    \begin{itemize}[nosep, leftmargin=0.5cm]
        \item PosX: Khu vực
        \item PosY: Số lượng listing
        \item Hue Color: Loại phòng
    \end{itemize}
\end{itemize}

\vspace{0.2cm}
\textbf{Manipulate:}
\begin{itemize}[nosep, leftmargin=*, topsep=0.1cm]
    \item Selection: Chọn một loại phòng để highlight trên toàn bộ biểu đồ.
    \item Hover để xem số lượng listing của từng loại phòng trong từng khu vực.
\end{itemize}

\vspace{0.2cm}
\textbf{Facet:}
\begin{itemize}[nosep, leftmargin=*, topsep=0.1cm]
    \item Có thể facet theo nhóm sức chứa nếu cần phân tích sâu hơn.
\end{itemize}

\vspace{0.2cm}
\textbf{Reduce:}
\begin{itemize}[nosep, leftmargin=*, topsep=0.1cm]
    \item Sắp xếp khu vực theo tổng số listing giảm dần.
\end{itemize} \\
\addlinespace[0.2cm]

\textbf{Scale} &
\begin{itemize}[nosep, leftmargin=*, topsep=0pt]
    \item Main key: khu vực
\end{itemize} \\
\end{longtable}

\begin{figure}[H]
    \centering
    \safeincludegraphics[width=0.9\linewidth]{charts/domain_task3_chart1.png}
    \caption{Hình 3.1. Biểu đồ cơ cấu loại phòng theo quận (Stacked Bar Chart)}
    \label{fig:domain-task3-chart1}
\end{figure}

\subsubsection*{B. Đánh giá biểu đồ}

\textbf{1. Nguyên lý biểu đạt (Expressiveness):}

\begin{itemize}
    \item Thuộc tính: Khu vực (Categorical -- C), Loại phòng (Categorical -- C), Số lượng listing (Quantitative -- Q).
    \item Channel:
    \begin{itemize}
        \item PosX: Định vị phân nhóm trên trục hoành $\rightarrow$ dùng cho thuộc tính Khu vực (C).
        \item PosY (Chiều dài/Length): Thể hiện độ lớn của giá trị định lượng $\rightarrow$ dùng cho thuộc tính Số lượng listing (Q).
        \item Color (Hue -- Sắc độ màu): Thể hiện sự phân loại danh mục $\rightarrow$ dùng cho thuộc tính Loại phòng (C).
    \end{itemize}
    \item Đánh giá mức độ quan trọng và phân bổ kênh theo Mackinlay:
    \begin{itemize}
        \item Ưu tiên 1 -- Chiều dài (PosY): Trong biểu đồ cột, yếu tố cốt lõi nhất là so sánh độ lớn. Theo thang đo Mackinlay, kênh Chiều dài (Length) trên một trục chung (aligned scale) là kênh biểu đạt mạnh nhất và chính xác nhất cho dữ liệu định lượng (Q). Việc gán ``Số lượng listing'' vào PosY đảm bảo tính chính xác tuyệt đối của biểu đồ.
        \item Ưu tiên 2 -- Vị trí không gian 1D (PosX): Quan trọng thứ hai để phân tách các khu vực địa lý (C) thành các nhóm rời rạc, tránh chồng lấp dữ liệu.
        \item Ưu tiên 3 -- Màu sắc (Color Hue): Thuộc tính Loại phòng (C) buộc phải dùng Color Hue. Theo Mackinlay, Color Hue là kênh biểu đạt hoàn hảo nhất cho dữ liệu định danh (Categorical), giúp mắt người phân loại tự nhiên mà không tạo ra ảo giác nhóm này ``ưu tiên'' hay ``lớn hơn'' nhóm kia.
    \end{itemize}
    \item Kết luận: Dùng hoàn toàn chuẩn xác các kênh biểu đạt. Cấu trúc biểu đồ Stacked Bar phản ánh đúng tính phân cấp và cộng dồn của dữ liệu.
\end{itemize}

\textbf{2. Nguyên lý hiệu quả (Effectiveness):}

\begin{itemize}
    \item Độ chính xác (Accuracy):
    \begin{itemize}
        \item Nhóm sub-bar dưới cùng (ví dụ: Private room màu đỏ) nằm trên một đường cơ sở chung (aligned baseline) (trục X = 0). Theo định luật Stevens' Psychophysical law, cảm nhận chiều dài 1D có độ lỗi cực kỳ thấp ($p \approx 1.0$), mức độ lỗi Log\_error gần như xấp xỉ $0$. Điều này giúp mắt người ước lượng phần base bar giữa các quận cực kỳ chính xác.
        \item Tuy nhiên, các sub-bar xếp chồng phía trên (Entire home/apt, Shared room) bị mất đường cơ sở chung (unaligned). Mắt người phải dùng kênh độ dài thuần túy thay vì đối chiếu vị trí điểm cuối, làm tăng sai số cảm nhận. Thao tác Manipulate (Hover xem Tooltip) được thiết lập nhằm bù đắp hoàn toàn nhược điểm này.
    \end{itemize}
    \item Khả năng phân biệt (Discriminability): Bảng màu Color Hue sử dụng 4 màu (Xanh dương, Cam, Đỏ, Xanh lá) có độ tương phản rất cao. Số lượng màu này nằm gọn trong ngưỡng ghi nhớ an toàn của mắt người (từ 6--12 sắc độ), giúp người xem lập tức nhìn ra các dải phân bổ loại phòng rõ ràng mà không bị nhầm lẫn hay quá tải thị giác.
    \item Khả năng tách biệt (Separability): Kênh vị trí (Pos) và kênh màu sắc (Color) tách biệt hoàn toàn (Fully separable) và không gây nhiễu chéo. Việc cột cao hay thấp không làm thay đổi cảm nhận về sắc độ màu, và ngược lại. Thao tác Reduce (sắp xếp khu vực theo tổng số listing giảm dần từ trái sang phải) dẫn dắt luồng mắt logic, giúp não bộ so sánh (Compare) trơn tru và giảm tải nhận thức (cognitive load) đáng kể.
\end{itemize}

\subsubsection*{C. Phân tích biểu đồ (Insight)}

\begin{itemize}
    \item Manhattan và Brooklyn thống trị hoàn toàn nguồn cung Airbnb, phản ánh sự tập trung nhu cầu lưu trú khổng lồ tại vùng lõi trung tâm New York.
    \item Thị trường Manhattan bị áp đảo tuyệt đối bởi ``Entire home/apt'', phục vụ tệp khách có ngân sách cao và cho thấy mức độ chuyên nghiệp hóa bất động sản lớn tại đây.
    \item Tại Brooklyn và Queens, loại hình ``Private room'' tăng vọt, chứng tỏ xu hướng người dân địa phương cho thuê không gian thừa để san sẻ gánh nặng chi phí sinh hoạt.
    \item Mô hình ``Shared room'' và ``Hotel room'' gần như vắng bóng, cho thấy xu hướng ở ghép đã thoái trào và khách sạn truyền thống chưa thực sự thâm nhập sâu vào nền tảng này.
\end{itemize}

\subsection{Biểu đồ 2: Sức chứa theo quận (Heat Map)}

\subsubsection*{A. Thiết kế Idiom}

\begin{longtable}{@{} p{2.8cm} p{12cm} @{}}
\toprule
\textbf{Đặc điểm} & \textbf{Chi tiết} \\
\midrule
\endhead

\bottomrule
\endfoot

\textbf{Idiom} & Heat Map (Bản đồ nhiệt) \\
\addlinespace[0.2cm]

\textbf{What} & Khu vực: Categorical, Nhóm sức chứa: Ordinal, Số lượng listing: Quantitative \\
\addlinespace[0.2cm]

\textbf{Why} & produce $\rightarrow$ browse, lookup $\rightarrow$ summarize
\begin{itemize}[nosep, leftmargin=*, topsep=0.1cm]
    \item Biểu đồ giúp xác định khu vực nào phù hợp với từng quy mô khách, ví dụ khu vực nào có nhiều listing cho cặp đôi, nhóm nhỏ, gia đình hoặc nhóm đông người.
\end{itemize} \\
\addlinespace[0.2cm]

\textbf{How} & \textbf{Encode:}
\begin{itemize}[nosep, leftmargin=*, topsep=0.1cm]
    \item Mark: Area
    \item Channel:
    \begin{itemize}[nosep, leftmargin=0.5cm]
        \item PosX: Nhóm sức chứa
        \item PosY: Khu vực
        \item Hue Color: Số lượng listing
    \end{itemize}
\end{itemize}

\vspace{0.2cm}
\textbf{Manipulate:}
\begin{itemize}[nosep, leftmargin=*, topsep=0.1cm]
    \item Selection: Chọn một ô để highlight tổ hợp khu vực và nhóm sức chứa.
    \item Hover để xem số lượng listing chi tiết.
\end{itemize}

\vspace{0.2cm}
\textbf{Facet:}
\begin{itemize}[nosep, leftmargin=*, topsep=0.1cm]
    \item Có thể facet theo loại phòng.
\end{itemize}

\vspace{0.2cm}
\textbf{Reduce:}
\begin{itemize}[nosep, leftmargin=*, topsep=0.1cm]
    \item Chuẩn hóa sức chứa thành các nhóm: 1--2 khách, 3--4 khách, 5--6 khách, 7+ khách.
\end{itemize} \\
\addlinespace[0.2cm]

\textbf{Scale} &
\begin{itemize}[nosep, leftmargin=*, topsep=0pt]
    \item Main key: khu vực $\times$ nhóm sức chứa
\end{itemize} \\
\end{longtable}

\begin{figure}[H]
    \centering
    \safeincludegraphics[width=0.9\linewidth]{charts/domain_task3_chart2.png}
    \caption{Hình 3.2. Biểu đồ sức chứa theo quận (Heat Map)}
    \label{fig:domain-task3-chart2}
\end{figure}

\subsubsection*{B. Đánh giá biểu đồ}

\textbf{1. Nguyên lý biểu đạt (Expressiveness):}

\begin{itemize}
    \item Thuộc tính: Khu vực (Categorical -- C), Nhóm sức chứa (Ordinal -- O), Số lượng listing (Quantitative -- Q).
    \item Channel:
    \begin{itemize}
        \item PosX: Nhóm sức chứa (Ordinal).
        \item PosY: Khu vực (Categorical).
        \item Color (Saturation -- Độ bão hòa màu): Thể hiện độ lớn của giá trị định lượng (Số lượng listing).
    \end{itemize}
    \item Đánh giá mức độ quan trọng và phân bổ kênh theo Mackinlay:
    \begin{itemize}
        \item Ưu tiên 1 -- Vị trí không gian (PosX, PosY): Sử dụng hai trục tọa độ để tạo thành lưới ma trận (Matrix) là phương án tối ưu để biểu đạt mối quan hệ giữa hai thuộc tính định danh/thứ tự. Theo Mackinlay, kênh Vị trí là kênh mạnh nhất cho mọi loại dữ liệu, giúp người xem dễ dàng tra cứu (Lookup) tổ hợp giữa Quận và Sức chứa.
        \item Ưu tiên 2 -- Màu sắc (Color Saturation/Luminance): Do hai kênh Vị trí đã bị chiếm dụng cho các chiều phân loại, giá trị định lượng (Q) phải sử dụng kênh màu sắc. Theo thang đo Mackinlay cho dữ liệu định lượng, kênh này xếp hạng thấp hơn Vị trí/Chiều dài nhưng lại cực kỳ hiệu quả trong việc nhận diện các ``điểm nóng'' (Hotspots) hoặc quy luật mật độ trên diện rộng.
    \end{itemize}
    \item Kết luận: Biểu đồ sử dụng chuẩn xác các kênh biểu đạt. Việc phân nhóm (Binning) sức chứa thành 4 cấp độ (1-2, 3-4, 5-6, 7+) giúp tinh gọn dữ liệu và tăng tính biểu đạt cho các User Story về quy mô khách.
\end{itemize}

\textbf{2. Nguyên lý hiệu quả (Effectiveness):}

\begin{itemize}
    \item Độ chính xác (Accuracy):
    \begin{itemize}
        \item Kênh vị trí (PosX, PosY) đảm bảo việc phân loại các nhóm sức chứa và quận huyện đạt độ chính xác tuyệt đối.
        \item Kênh màu sắc (Color Saturation) tuân theo định luật Stevens' Psychophysical law với hàm mũ cảm nhận thấp ($p < 1.0$), dẫn đến sai lệch cảm nhận về lượng (ví dụ: khó phân biệt một ô đậm hơn 10\% có thực sự nhiều hơn 10\% số lượng hay không). Độ lỗi Log\_error ước tính rơi vào khoảng $\approx 2.5$. Tuy nhiên, thiết kế đã khắc phục hoàn toàn nhược điểm này bằng cách hiển thị trực tiếp nhãn số (Labels) bên trong các ô vuông (Mark Area).
    \end{itemize}
    \item Khả năng phân biệt (Discriminability): Dải màu chuyển sắc (Sequential) giúp phân biệt rõ các vùng có mật độ cao (Manhattan 1-2 khách: 5,541) so với các vùng khan hiếm (Staten Island 7+ khách: 16). Khoảng cách giữa các ô (Grid) giúp tách biệt rõ ràng các đơn vị dữ liệu, tránh hiện tượng hòa lẫn màu sắc.
    \item Khả năng tách biệt (Separability): Các kênh vị trí và màu sắc tách biệt tốt. Thao tác Reduce (Chuẩn hóa sức chứa thành các nhóm Ordinal) giúp loại bỏ sự nhiễu loạn của hàng chục giá trị số lẻ, tập trung sự chú ý của người xem vào các phân khúc thị trường chính.
\end{itemize}

\subsubsection*{C. Phân tích biểu đồ (Insight)}

\begin{itemize}
    \item Phân khúc phòng dành cho 1-2 khách (khách lẻ hoặc cặp đôi) chiếm đa số tuyệt đối tại mọi khu vực, đặc biệt bùng nổ tại Manhattan và Brooklyn.
    \item Nguồn cung tỷ lệ nghịch với sức chứa: số lượng listing giảm mạnh khi quy mô nhóm khách tăng lên, cho thấy Airbnb NYC chủ yếu phục vụ tệp khách nhỏ.
    \item Manhattan và Brooklyn là hai ``điểm nóng'' cung cấp giải pháp lưu trú cho mọi quy mô nhóm khách, trong khi Staten Island gần như đứng ngoài phân khúc nhóm đông người (7+ khách).
    \item Sự tập trung mật độ cao tại ô ``Manhattan / 1-2 khách'' (5,541) xác nhận đây là phân khúc cạnh tranh khốc liệt nhất nhưng cũng có nhu cầu cao nhất toàn thị trường.
\end{itemize}
